% --- Archon physics formalization source begin ---
% archon:chemistry
% archon:covers IChO2026Problems/problem_icho_2026_t8_a6.lean
% archon:problem-id icho_2026_t8
% archon:part-id T8-A6
% archon:previous-part icho_2026_t8_a5
% archon:previous-part-policy derive_in_answer_blind_run_or_use_problem_stated_fallback

\chapter{Icho Chemistry problem icho\_2026\_t8\_a6}
\label{ch:IChO2026Problems_problem_icho_2026_t8_a6}

\paragraph{Problem source.}
\#\# Shared official problem context
T8. Recycling of Carbon Dioxide 7\% Removing CO2 gas from the atmosphere is key to combat climate change. One promising solution is the photocatalytic reduction of CO2 to CO using molecular catalyst 1. 8.1 Write the half equation for the reduction of CO2 to CO in an acidic medium. The conjugate oxidation reaction occurs with sacrificial reductant 2 in aqueous solution with the following mechanism: Species 3, 4, and 5 are short-lived ionic and radical intermediates. 7 yields a silver mirror with the [Ag(NH3)2]OH test. 8.2 Draw the structures of 3-7. Catalyst 1 can be synthesised with the following reaction. 8.3 Tick the correct geometry for the structure of 1. When 1 is loaded on crystalline carbon nitride (C3N4) support material, it can reduce CO2 via the following mechanism. 8.4 Draw the structures of 9-15. Use the cartoon structure for ligand 8 shown above. For 9-15, specify the oxidation state, 𝑂𝑆, coordination number, 𝐶𝑁, and number of valence electrons, 𝑉𝐸, of Fe, along with the total charge of corresponding complex structures, 𝑍. Note: Some information is provided on the answer sheet. 8.5 Calculate the number of catalytic molecules per nm2, 𝑁𝑐𝑎𝑡, of C3N4 loaded with 1, when the mass fraction of catalyst, 𝜔𝑐𝑎𝑡= 3.8\%, if the specific surface area of C3N4 is 17.8 m2 g−1. M𝑐𝑎𝑡= 557.21 g mol−1. The Turnover Frequency of a catalyst (TOF) is the number of product molecules formed per active catalyst molecule per hour. Under LED illumination (𝜆= 390 nm) with power of 50 mW, there was a TOF of 8 h−1 for 10 mg of C3N4, loaded with 1 with 𝜔𝑐𝑎𝑡= 3.8\% The formula for calculating the quantum yield is given below: 𝜑= 𝑛𝑢𝑚𝑏𝑒𝑟𝑜𝑓𝑟𝑒𝑎𝑐𝑡𝑒𝑑𝑒𝑙𝑒𝑐𝑡𝑟𝑜𝑛𝑠 𝑛𝑢𝑚𝑏𝑒𝑟𝑜𝑓𝑖𝑛𝑐𝑖𝑑𝑒𝑛𝑡𝑝ℎ𝑜𝑡𝑜𝑛𝑠⋅100\%

\#\# Current subquestion T8-A6
Calculate the quantum yield, 𝜑(\%), for CO formation.

\paragraph{Shared problem context.}
T8. Recycling of Carbon Dioxide 7\% Removing CO2 gas from the atmosphere is key to combat climate change. One promising solution is the photocatalytic reduction of CO2 to CO using molecular catalyst 1. 8.1 Write the half equation for the reduction of CO2 to CO in an acidic medium. The conjugate oxidation reaction occurs with sacrificial reductant 2 in aqueous solution with the following mechanism: Species 3, 4, and 5 are short-lived ionic and radical intermediates. 7 yields a silver mirror with the [Ag(NH3)2]OH test. 8.2 Draw the structures of 3-7. Catalyst 1 can be synthesised with the following reaction. 8.3 Tick the correct geometry for the structure of 1. When 1 is loaded on crystalline carbon nitride (C3N4) support material, it can reduce CO2 via the following mechanism. 8.4 Draw the structures of 9-15. Use the cartoon structure for ligand 8 shown above. For 9-15, specify the oxidation state, 𝑂𝑆, coordination number, 𝐶𝑁, and number of valence electrons, 𝑉𝐸, of Fe, along with the total charge of corresponding complex structures, 𝑍. Note: Some information is provided on the answer sheet. 8.5 Calculate the number of catalytic molecules per nm2, 𝑁𝑐𝑎𝑡, of C3N4 loaded with 1, when the mass fraction of catalyst, 𝜔𝑐𝑎𝑡= 3.8\%, if the specific surface area of C3N4 is 17.8 m2 g−1. M𝑐𝑎𝑡= 557.21 g mol−1. The Turnover Frequency of a catalyst (TOF) is the number of product molecules formed per active catalyst molecule per hour. Under LED illumination (𝜆= 390 nm) with power of 50 mW, there was a TOF of 8 h−1 for 10 mg of C3N4, loaded with 1 with 𝜔𝑐𝑎𝑡= 3.8\% The formula for calculating the quantum yield is given below: 𝜑= 𝑛𝑢𝑚𝑏𝑒𝑟𝑜𝑓𝑟𝑒𝑎𝑐𝑡𝑒𝑑𝑒𝑙𝑒𝑐𝑡𝑟𝑜𝑛𝑠 𝑛𝑢𝑚𝑏𝑒𝑟𝑜𝑓𝑖𝑛𝑐𝑖𝑑𝑒𝑛𝑡𝑝ℎ𝑜𝑡𝑜𝑛𝑠⋅100\%

\paragraph{Current subquestion.}
Calculate the quantum yield, 𝜑(\%), for CO formation.

\paragraph{Figure/image paths.}
\begin{itemize}
\item icho\_2026\_source/image/T8\_page-2.png
\item icho\_2026\_source/image/T8\_page-1.png
\end{itemize}

\paragraph{Reusable previous-part conclusions.}
\begin{itemize}
\item Source T8-A5. Question: Calculate the number of catalytic molecules per nm2, 𝑁𝑐𝑎𝑡, of C3N4 loaded with 1, when the mass fraction of catalyst, 𝜔𝑐𝑎𝑡= 3.8\%, if the specific surface area of C3N4 is 17.8 m2 g−1. M𝑐𝑎𝑡= 557.21 g mol−1. Policy: derive\_in\_answer\_blind\_run\_or\_use\_problem\_stated\_fallback
\end{itemize}

\paragraph{Formalization target.}
create a compiling Lean file with sorry bodies at `IChO2026Problems/problem\_icho\_2026\_t8\_a6.lean`.
The Lean declarations must preserve every mathematical object, quantifier, hypothesis, side condition, approximation/error guarantee, and requested conclusion from the source statement.
Use Mathlib/Physlib/CRNT names found through LeanExplore where available. Prefer the configured benchmark Base library; introduce a local definition only when the source genuinely requires an object that the environment does not expose.
This is an answer-blind solve-phase task. Derive a candidate result only from the problem-side evidence above; do not assume a target value, select a tolerance around a desired value, or encode a desired result into a candidate domain.
For numerical questions, define the raw end-to-end quantity and apply an explicit source-derived rounding rule. For identification questions, characterize or prove uniqueness before recording the derived candidate.

\begin{theorem}[Icho Chemistry formalization target]
\label{thm:physics:icho_2026_t8_a6:target}
The assigned autoformalize agent should translate this IChO chemistry problem into Lean declarations in the covered file, with theorem and lemma proof bodies written as `by sorry`.
\end{theorem}
\begin{proof}
This is an autoformalization task, not a proof task. Produce faithful statements that can later be proved without weakening the source contract.
\end{proof}
% --- Archon physics formalization source end ---
